% Created 2018-04-02 Mon 04:55
% Intended LaTeX compiler: pdflatex
\documentclass[final]{beamer}
	         \usetheme{ph}
           \usepackage[orientation=portrait,size=a0,scale=1.4]{beamerposter}
           \usepackage[absolute,overlay]{textpos}
	         \usepackage[authoryear]{natbib}


\setlength{\paperwidth}{36in}
\setlength{\paperheight}{48in}
\setlength{\textwidth}{0.98\paperwidth}
\setlength{\textheight}{0.98\paperheight}
\graphicspath{{../output/figures/}{../lib/}}
\usepackage[export]{adjustbox}
\usepackage{graphicx,caption}
\usepackage{eurosym}
\usepackage{textcomp}
\usepackage{bibentry}
\date{}
\author{
Philipp Homan$^{1}$,
\\
\normalsize{$^{1}$Department of Psychiatry,}
\normalsize{The Donald and Barbara Zucker}
\normalsize{School of Medicine at Northwell/Hofstra,}
\normalsize{Hempstead, NY}\\
}
\usetheme{default}
\date{2018-04-02 04:55}
\title{A scientific poster entirely written in org-mode using GNU emacs and the beamer library}
\begin{document}

\begin{frame}[label={sec:org8f3d3e7}]{}
\begin{columns}
\begin{column}[t]{0.45\textwidth\columnwidth}
\begin{block}{Background}
\begin{itemize}
\item Org-mode is not only useful for producing blog posts and even
scientific manuscripts; it is also perfectly suitable to produce
decent looking scientific posters.
\item We combine a relatively simple custom \LaTeX style file and common
org-mode syntax
\end{itemize}
\end{block}

\begin{block}{Methods: Part 1}
\captionsetup{justification=justified,width=.85\linewidth}
\begin{figure}[htbp]
\centering
\includegraphics[width=0.85\textwidth]{org-mode-poster-1.png}
\caption{\label{fig:orga7530f0}
\textbf{This could be a figure.}}
\end{figure}
\end{block}

\begin{block}{Methods: Part 2}
\captionsetup{justification=justified,width=.85\linewidth}
\begin{figure}[htbp]
\centering
\includegraphics[width=0.85\textwidth]{org-mode-poster-2.png}
\caption{\label{fig:org95c7011}
\textbf{This could be another figure.}}
\end{figure}
\end{block}
\end{column}

\begin{column}[t]{0.45\textwidth\columnwidth}
\begin{block}{Results: Part 1}
\captionsetup{justification=justified,width=.8\linewidth}
\begin{figure}[htbp]
\centering
\includegraphics[width=0.8\textwidth]{org-mode-poster-3.png}
\caption{\label{fig:org89b785d}
\textbf{Another figure}}
\end{figure}
\end{block}

\begin{block}{Results: Part 2}
\begin{columns}
\begin{column}[T]{0.48\columnwidth}
\captionsetup{justification=justified,width=.8\linewidth}
\begin{figure}[htbp]
\centering
\includegraphics[page=3,width=0.9\textwidth]{org-mode-poster-4.png}
\caption{\label{fig:org51482b3}
\textbf{This is another figure}}
\end{figure}
\end{column}

\begin{column}[T]{0.48\columnwidth}
\captionsetup{justification=justified,width=.8\linewidth}
\begin{figure}[htbp]
\centering
\includegraphics[page=9,width=0.9\textwidth]{org-mode-poster-4.png}
\caption{\label{fig:org2cc1106}
\textbf{This is the right column of this box}}
\end{figure}
\end{column}
\end{columns}
\end{block}

\begin{block}{Conclusions}
\begin{itemize}
\item This example shows how incredibly versatile org-mode is. One can now
produce scientific posters with a text editor (preferably emacs).
\end{itemize}
\end{block}
\end{column}
\end{columns}
\end{frame}
\end{document}
